% Shared module: the worked example (one complete run).
\subsection{A complete run on one data set}\label{sec:worked}
We follow one input through every stage. The data are the running example (\cref{ex:tpd}) with $c=1$,
$\nu=0.3$, observed on a $64\times128$ grid over $[0,2]\times[0,2\pi)$ with white noise of standard deviation
$\noise=0.05\,\mathrm{std}\abs u$, the hardest grid setting of the benchmark in \cref{sec:bench}. The pipeline is
run exactly as in the benchmark: one noisy grid, the outer loop with default parameters, and $16$
parametric-bootstrap refits for the intervals. The run takes about two minutes on one core.
The figures of this subsection also use a second run on the same input with $\Nsamp=32$ posterior samples,
which shows the cloud of \cref{sec:bayes-cloud}.

\paragraph{The band} Of the 64 positive wavenumbers below $\pi/\Delta x=64$, only $\kappa\le4$ rise above the
detection threshold (\cref{fig:nyquist}); the noise-free field has negligible energy beyond $\kappa\approx5$, and
a diffusive mode at $\kappa=3$ loses a factor $e^{-0.3\cdot9\cdot2}\approx e^{-5.4}$ of its amplitude over the
window. Every later decision is made with these few bins.


\paragraph{Modes and directions} The matrix pencil returns twelve modes (\cref{fig:directions}). Eight lie on
the line $\operatorname{Im}\lambda=-\kappa$ with $\operatorname{Re}\lambda\approx0$ (transport), four on the parabola
$\operatorname{Re}\lambda=-0.3\kappa^2$ with $\operatorname{Im}\lambda\approx0$ (diffusion). Each factor is a
\emph{direction} in the frequency plane: a curve along which its modes line up. The structured fit
recovers the two curves; the thin lines are the factors fitted to the 32 posterior samples and show the
uncertainty of each curve, larger for diffusion at high $\abs\kappa$, where there are no modes. The counting
condition of \cref{prop:counting} holds: the transport modes occupy eight distinct wavenumbers
($\kappa=\pm1,\dots,\pm4$) and the diffusion modes four ($\pm1,\pm2$), against the two and three required
for factors of degree one and two in $\xi$.

\emph{How \cref{fig:directions} is built.} Each point is one mode $(\kappa,\lambda)$ returned by the pencil: its
horizontal position is the wavenumber, its height the oscillation $\operatorname{Im}\lambda$ (panel a) or the damping
$\operatorname{Re}\lambda$ (panel b), and its colour the factor to which the structured fit assigned it. Each thin line is
one fitted factor, drawn as its dispersion curve $\lambda=\Lambda_j(i\kappa)$ for all $\kappa$, one line per factor and
posterior sample. In the figure, the transport group is a sloped line in (a) and a flat line at zero in (b), a mode
that travels without decaying; the diffusion group is flat in (a) and a downward parabola in (b), a mode that does not
travel and decays faster at higher wavenumber. The thin lines bundle tightly where there are modes and fan out where
there are none.

\begin{figure}[tbp]
\centering
\includegraphics[width=\linewidth]{fig_directions}
\caption{The two groups of the running example as directions in the frequency plane. Points: modes found
by the matrix pencil in the noisy data, coloured by the factor they were assigned to. Dashed: the true
dispersion branches $\lambda=-ic\kappa$ and $\lambda=-\nu\kappa^2$. Thin lines: the factors fitted to 32
posterior samples of the same input; their spread is the uncertainty of each direction. The points come from the run
with the benchmark protocol, the lines from the 32-sample run. (a) Oscillation, $\operatorname{Im}\lambda$;
(b) damping, $\operatorname{Re}\lambda$.}\label{fig:directions}
\end{figure}

\paragraph{Grouping and the evidence step} At the starting class the data-level evidence ranks the
structure $((1,1),(1,2))$ (one factor of degree $(1,1)$, one of degree $(1,2)$) first, $0.13$ nats ahead of
$((1,1),(1,3))$; the second class contains the first, so both are compatible with the truth and the simpler
one is kept. The core, applied to the noisy grid, initially returned \emph{three} factors: the diffusion
factor and two near-copies of the transport factor, $\tau+1.0006\,\xi$ and $\tau+1.0044\,\xi-0.0028$. Grouping
by evidence on the observations removed the duplicate (the factor set without it has higher evidence), and
dropped the constant terms, whose intervals contained zero. An earlier version of the pipeline (the row ``previous version'' of \cref{tab:bench}) returned a similar triple on
this input, two transport copies $\tau+0.9989\,\xi$ and $\tau+1.0044\,\xi$ and a diffusion factor $\tau-0.263\,\xi^2$, and
labelled it \FC: a confident wrong answer that the grouping step now prevents.

\paragraph{Cloud and clusters} Over 32 posterior samples the core returns 64 factors. In coefficient space
they form two tight clusters around the true factors (\cref{fig:grouping}a), and the largest principal angle
between their kernels is near zero inside a cluster and about $1.2$ rad across clusters
(\cref{fig:grouping}b). The Dirichlet-process mixture finds two clusters with masses $0.767$ and $0.233$;
$P(k=2)=1$ and $\keff=1.72$ ($90\%$ interval $[1.67,1.79]$), which is far from $2$ only because the two
groups carry unequal energy.

\emph{How \cref{fig:grouping} is built.} In panel (a) each marker is one factor found in one posterior sample,
placed at its coefficient of $\xi$ (horizontal) and of $\xi^2$ (vertical), with area proportional to its mass in the
cloud. Panel (b) compares every pair of these factors by the spaces of functions they allow: for each factor we take the
span of its kernel exponentials on the grid and band, and compute the principal angles between the two spans from the
singular values of $B_i^HB_j$ (their cosines). The colour shows the largest of these angles, light for nearly
the same space and dark for very different spaces, with rows and columns ordered by cluster. The figure shows, in (a),
two tight, well separated clouds and nothing in between; in (b) two light blocks on the diagonal, because all factors of
one group allow the same functions, and dark blocks off the diagonal, because the two groups do not. (The largest angle
is used here to tell kernels apart; the Friedrichs angle of \cref{sec:cond} is the smallest nonzero one, which measures
how well two different groups can be separated.)

\begin{figure}[tbp]
\centering
\includegraphics[width=\linewidth]{fig_grouping}
\caption{Grouping across posterior samples. (a) Each marker is one factor returned for one of 32 posterior
samples, placed by its coefficients of $\xi$ and $\xi^2$, the two that separate the groups; area is proportional to its mass $\pi_qw_j$; crosses mark the true
factors. (b) Largest principal angle between the sampled kernels of every pair of factors, ordered by cluster:
the block structure is the grouping.}\label{fig:grouping}
\end{figure}

\paragraph{Components and energy} \Cref{fig:components} shows the data, the two components and the residual.
The transport group carries $w_1=0.768$ of the energy and the diffusion group $w_2=0.232$, with $90\%$
intervals of width below $0.002$. The joint residual of the data on the two kernels is $2.5\%$ of the data
norm, and it shows no structure (\cref{fig:components}d). Each reported factor has a single term besides $\partial_t$, so its
interaction ratio is $\rho=0$ by definition. Computed on the factors of the cloud before pruning, $\rho$ has median $0.03$
($90\%$ interval $[0.01,0.07]$ over 32 samples), entirely due to the small constant and $\xi$ coefficients that
the evidence step removes. For comparison, on the advection--diffusion equation (one factor with the two
terms $c\partial_x$ and $\nu\partial_x^2$) the same pipeline reports $\rho=0.79$, with a $90\%$ interval
$[0.788,0.793]$ over 16 samples: there the interaction is most of the evolution, as \cref{thm:eint} predicts.

\emph{How \cref{fig:components} is built.} Each panel is a map of a function of $(t,x)$, with $x$ horizontal and $t$
vertical, red for positive and blue for negative values, on one colour scale for all panels. The components are the sums
of the modes of each group, \eqref{eq:uj}. The transport component consists of tilted stripes, a profile
that moves at speed $c=1$ without changing shape; the diffusion component has vertical structure that fades with time; the
residual has no pattern left.

\begin{figure}[tbp]
\centering
\includegraphics[width=\linewidth]{fig_components}
\caption{The data of the worked example and its two groups as functions of $(t,x)$, on a common colour scale:
(a) noisy data; (b) transport component; (c) diffusion component; (d) residual, magnified five times, which
is indistinguishable from the noise.}\label{fig:components}
\end{figure}

\paragraph{The output} \Cref{tab:worked} lists the final output and how to read it. The separation of the two
groups is moderate: $\cF=0.827$, $\thetaF=0.60$ rad, so errors in the components are amplified by at most
$\gN^{-1/2}=2.41$. For contrast, the two conjugate factors of the beam equation
$\tau\mp0.3\,i\xi^2$ (\cref{app:cases}) are recovered with the same accuracy but are much closer to each other: $\cF=0.939$, an
amplification of $4.05$.

\begin{table}[tbp]
\centering\small
\caption{Output of the worked example and how to read each quantity. Intervals are $90\%$ intervals
from 16 parametric-bootstrap refits.}\label{tab:worked}
\begin{tabularx}{\linewidth}{@{}l l X@{}}
\toprule
Quantity & Value & Meaning\\
\midrule
equation & $(\tau+1.0006\,\xi)(\tau-0.2980\,\xi^2)$ & recovered factors; truth $(\tau+\xi)(\tau-0.3\,\xi^2)$, largest relative error $0.7\%$\\
coefficient of $\xi$ & $1.0006\ [0.9998,1.0020]$ & speed $c$ of the transport group\\
coefficient of $\xi^2$ & $-0.2980\ [-0.3049,-0.2919]$ & diffusivity $\nu$ of the diffusion group\\
label & \FC & enough distinct wavenumbers per factor, goodness of fit passed\\
$P(k)$ & $P(k=2)=1$ & posterior over the number of groups\\
$w_1,\ w_2$ & $0.768,\ 0.232$ & energy shares~\eqref{eq:energy} of the two groups\\
$P(\text{split})$ & $1.000$ & evidence that the data need both groups\\
$\gN,\ \cF,\ \thetaF$ & $0.173,\ 0.827,\ 0.60$ & separation of the groups (\cref{sec:cond})\\
$\gN^{-1/2}$ & $2.41$ & worst-case error amplification of the components\\
$\rho$ & $0$ & no interaction inside either reported factor (single terms)\\
\bottomrule
\end{tabularx}
\end{table}
